\appendix
\chapter{Test Suite Evidence}
\label{app:test-suite-evidence}

\section{Complete Test File Inventory}
\label{app:test-inventory}

Table~\ref{tab:test-inventory} lists all 15 test files in the \texttt{wasm4pm-cognition} crate,
with exact test counts, oracle rank, and coverage area as verified on 2026-06-09.

\begin{table}[ht]
\centering
\small
\begin{tabular}{llrl}
\hline
\textbf{File} & \textbf{Tests} & \textbf{Oracle Rank} & \textbf{Coverage Area} \\
\hline
\texttt{tests/breed\_adversarial.rs}             & 58 & L1 & Known-input oracles + bypass attempts, all 13 breeds \\
\texttt{tests/breed\_oracle\_gaps.rs}            & 31 & L1 & Mathematical property invariants \\
\texttt{tests/autoinstinct\_adversarial.rs}      & 27 & L1 & Unfakeable autoinstinct oracles \\
\texttt{tests/strips\_soar\_cbr\_invariants.rs}  & 23 & L1 & Formal invariants (STRIPS, SOAR, CBR) \\
\texttt{tests/dispatch\_smoke.rs}                & 20 & L2 & Routing smoke tests (all breeds reachable) \\
\texttt{tests/gps\_dendral\_eliza\_falsification.rs} & 19 & L1 & Formal property tests (GPS, DENDRAL, ELIZA) \\
\texttt{tests/breed\_determinism.rs}             & 16 & L1 & 13-breed bit-exact double-run harness \\
\texttt{tests/breed\_math\_properties.rs}        & 12 & L1 & Property-based CF and similarity tests \\
\texttt{tests/level\_10\_integration.rs}         & 11 & L3 & End-to-end integration (full pipeline) \\
\texttt{tests/strips\_gps\_level\_10.rs}         & 10 & L3 & STRIPS/GPS integration, plan soundness \\
\texttt{tests/adversarial\_bypass.rs}            & 10 & L1 & Bypass/injection hardening \\
\texttt{tests/ocel\_conformance.rs}              &  7 & L1 & OCEL lifecycle conformance enforcement \\
\texttt{tests/combine\_cf\_properties.rs}        &  5 & L1 & Proptest CF formula (Shortliffe-Buchanan) \\
\texttt{tests/authority\_classifier\_property.rs}&  4 & L1 & Authority classifier invariants \\
\texttt{tests/ghf\_fleet\_sentinel\_test.rs}     &  1 & L3 & GHF proof gate sentinel \\
\hline
\textbf{Total}                                   & \textbf{254} & & \textbf{wasm4pm-cognition Rust tests} \\
\hline
\end{tabular}
\caption{Complete test file inventory for \texttt{wasm4pm-cognition} (verified 2026-06-09).
Oracle rank L1 = known-input or property-based unfakeable oracle;
L2 = routing/dispatch coverage;
L3 = integration/end-to-end.}
\label{tab:test-inventory}
\end{table}

\noindent
System-wide totals (2026-06-09): 480 passing Rust tests, 175 passing TypeScript tests,
655 total across the system. Zero failures.

\section{The 13 Lifecycle Models}
\label{app:lifecycle-models}

Each \texttt{BreedLifecycleModel} is a constant in \texttt{src/ocel.rs} that encodes the
lawful phase sequence for one breed as a deterministic finite automaton.
The pseudo-regex notation below uses \texttt{*} (zero or more), \texttt{+} (one or more),
\texttt{?} (zero or one), and \texttt{|} (alternation).
All 13 constants are enforced at the \texttt{run\_breed} choke point in \texttt{wasm.rs}.

\begin{description}

\item[MYCIN]
\begin{verbatim}
load-fact* -> fire-rule+ -> decision?
\end{verbatim}
Facts are loaded (possibly zero), at least one rule fires, and the run optionally
reaches a decision threshold.

\item[Hearsay]
\begin{verbatim}
(post-hypothesis | enqueue-ksar | stale-ksar)+ -> decision?
\end{verbatim}
The blackboard cycle is opportunistic: hypotheses, KSAR enqueue, and stale-KSAR
re-rating may interleave freely; at least one event must occur.

\item[CBR]
\begin{verbatim}
build-index -> retrieve-candidates
           -> (score-case | reuse-adapt | revise-* | retain-case)*
           -> decision?
\end{verbatim}
The index must be built before retrieval; the four-R adaptation cycle is optional
per case; decision is optional if the case base is empty.

\item[GPS]
\begin{verbatim}
(match-goal | set-goal)+
-> (apply-operator | subgoal | achieve-diff)*
-> (decision | no-plan)?
\end{verbatim}
At least one goal event must be recorded; operator application, subgoaling,
and difference reduction are optional; the run ends with a decision or
explicit \texttt{no-plan}.

\item[STRIPS]
\begin{verbatim}
(apply-action | subgoal | add-effect | del-effect)+
-> (decision | no-plan)?
\end{verbatim}
At least one action or effect event must occur; the run terminates with a
decision or \texttt{no-plan}.

\item[Prolog]
\begin{verbatim}
load-fact* -> (unify | sld-step | match-rule | bind-var)+ -> decision?
\end{verbatim}
Facts are loaded, then at least one resolution step occurs, ending
optionally with a decision (bound answer or failure).

\item[SOAR]
\begin{verbatim}
(propose-operator | preference)+
-> (decide-operator | impasse | subgoal)+
-> (apply-operator | decision)*
\end{verbatim}
Proposals and preferences must precede any decision cycle; an impasse
triggers a subgoal; application and decisions may repeat.

\item[ELIZA]
\begin{verbatim}
(match-pattern | keyword-match)+
-> (transform | reflect | response | decision)*
\end{verbatim}
At least one pattern or keyword event; the transformation/response
pipeline is optional (degenerate input may produce no output).

\item[DENDRAL]
\begin{verbatim}
(generate-hypothesis | enumerate)+
-> (test-hypothesis | eliminate | decision)*
\end{verbatim}
Hypothesis generation must precede any testing; elimination and decision
are optional (an exhausted search space reaches \texttt{decision} with
an empty candidate set).

\item[AutoinstinctNeurosis]
\begin{verbatim}
(analyze | detect-pattern | belief-update)+ -> decision?
\end{verbatim}
Iterative analysis and belief updates; decision is optional if no
anomalous belief exceeds the threshold.

\item[AutoinstinctVision]
\begin{verbatim}
(perceive | segment | classify)+ -> decision?
\end{verbatim}
Perception and segmentation must precede classification; at least one
event required; decision is optional.

\item[AutoinstinctSemantics]
\begin{verbatim}
(parse | frame-bind | atrans | ptrans)+ -> decision?
\end{verbatim}
Parsing and frame binding interleave with conceptual-dependency transfers
(\textsc{atrans}/\textsc{ptrans}); decision is optional.

\item[AutoinstinctLearning]
\begin{verbatim}
(plan-step | update-distance | expand-frontier)+
-> (goal-reached | decision)?
\end{verbatim}
Each planning step updates the heuristic distance; frontier expansion
is optional; the run terminates with \texttt{goal-reached} or a
\texttt{decision}; monotone distance non-increase is a structural invariant.

\end{description}

\section{The Prolog Grandparent Test (Annotated)}
\label{app:prolog-grandparent}

The test below is reproduced from \texttt{crates/prolog8/} and proves that
Robinson flat-term unification with shared-variable propagation is correctly
implemented.  A stub or regex-based mock cannot pass this test because it
requires binding \texttt{?1} (the intermediate) and propagating that binding
to both occurrences of the same variable slot.

\begin{verbatim}
#[test]
fn test_grandparent_unification() {
    // Build a minimal knowledge base:
    //   parent(alice, bob).
    //   parent(bob,   carol).
    //   grandparent(?0, ?2) :- parent(?0, ?1), parent(?1, ?2).
    //
    // ?N notation = positional (De-Bruijn-like) variable slots.
    // The engine must NOT treat ?1 as two independent fresh variables --
    // the same binding must flow from the first parent/2 call into the
    // second.  This is the minimal test of shared-variable unification.
    let mut kb = KnowledgeBase::new();

    // EDB facts: parent/2
    kb.assert_fact("parent", &[Term::atom("alice"), Term::atom("bob")]);
    kb.assert_fact("parent", &[Term::atom("bob"),   Term::atom("carol")]);

    // IDB rule: grandparent(?0,?2) :- parent(?0,?1), parent(?1,?2)
    // ?0 = first arg (grandparent's ancestor)
    // ?1 = shared intermediate (must unify consistently across both body literals)
    // ?2 = second arg (grandparent's descendant)
    kb.assert_rule(
        Clause::rule(
            Atom::new("grandparent", vec![Term::var(0), Term::var(2)]),
            vec![
                Atom::new("parent", vec![Term::var(0), Term::var(1)]),
                Atom::new("parent", vec![Term::var(1), Term::var(2)]),
            ],
        )
    );

    // Query: grandparent(alice, ?X)
    // Expected answer: ?X = carol
    //
    // SLD trace:
    //   1. Resolve grandparent(alice,?X) against the rule head:
    //      unify(?0=alice, ?2=?X) -> subst: {?0->alice, ?2->?X}
    //   2. First body literal: parent(alice, ?1)
    //      unify with fact parent(alice,bob) -> subst: {?1->bob}
    //      CRITICAL: ?1 is now locked to "bob" for the entire resolvent.
    //   3. Second body literal: parent(bob, ?X)
    //      (substituting ?1->bob into parent(?1,?2) gives parent(bob,?2))
    //      unify with fact parent(bob,carol) -> subst: {?X->carol}
    //   4. Empty resolvent -> success.  Answer: carol.
    //
    // If ?1 were treated as two independent fresh variables, step 3 would
    // attempt parent(<fresh>,?X) and match parent(alice,bob), yielding
    // ?X=bob (wrong) or match parent(bob,carol) accidentally (coincidence,
    // not proof).  Only correct shared-variable unification yields carol
    // deterministically.
    let engine = SldEngine::new(&kb);
    let query  = Atom::new("grandparent", vec![Term::atom("alice"), Term::var(0)]);
    let results = engine.solve(&query, /*depth_limit=*/10);

    assert_eq!(results.len(), 1,
        "Expected exactly one proof path (alice->bob->carol)");
    assert_eq!(
        results[0].get_binding(0),   // ?X bound in the answer substitution
        Some(&Term::atom("carol")),
        "Shared-variable unification must propagate ?1=bob into the \
         second parent/2 call and yield carol, not any other atom"
    );
}
\end{verbatim}

\noindent
\textbf{Why this is unfakeable.}
The test requires (1) a full unification engine that propagates bindings across
shared variables, (2) SLD resolution that consults multiple body literals in
order, and (3) correct answer-substitution extraction.  No look-up table or
string match can satisfy the assertion because the answer (\textit{carol}) is
derived, not stored directly under any query term.

\section{The OCEL Conformance Refusal Mechanism}
\label{app:ocel-refusal}

The listing below is the \texttt{run\_breed} choke point in
\texttt{crates/wasm4pm-cognition/src/wasm.rs} with the OCEL enforcement path
annotated.  Non-conforming execution is refused before any receipt is emitted.

\begin{verbatim}
// wasm.rs -- run_breed (simplified for readability; annotations in comments)

pub fn run_breed(input: BreedInput) -> Result<ContractResult, CognitionError> {

    // -- 1. Dispatch to breed implementation ----------------------------------
    // Each breed returns (BreedOutput, Vec<TraceStep>).
    // TraceStep records (phase_label, logical_step: u64) -- no wall-clock.
    let (output, trace) = dispatch_breed(&input)?;

    // -- 2. Derive OCEL log from trace -----------------------------------------
    // derive_ocel converts raw TraceSteps into an OcelLog with:
    //   * one OcelObject per run (object_type = breed name)
    //   * one OcelEvent per trace step (activity = phase_label)
    //   * ocel:timestamp = logical_step (integer, strictly monotone)
    //
    // No wall-clock is used -- OCEL timestamps are logical sequence numbers.
    // This makes the log deterministic and replay-safe.
    let ocel_log = derive_ocel(&input.breed, &output.run_id, &trace);

    // -- 3. Select the lifecycle model for this breed --------------------------
    // Each breed has a BreedLifecycleModel constant (see Appendix A.2).
    // The model is a DFA over phase labels.
    let model = lifecycle_model_for(&input.breed)
        .ok_or(CognitionError::UnknownBreed(input.breed.clone()))?;

    // -- 4. Validate OCEL log against lifecycle model --------------------------
    // validate_ocel_alignment replays the OcelLog through the DFA.
    // Returns ConformanceResult { fitness, is_conforming, violations }.
    //
    // REFUSAL POINT: if the log does not conform, the entire run is aborted.
    // The caller receives an Err -- no receipt is emitted, no output_hash
    // is written, and no run_id is recorded.  The process event log cannot
    // lie about what happened because the lie would be detected here.
    let conformance = validate_ocel_alignment(&ocel_log, &model)?;
    if !conformance.is_conforming {
        return Err(CognitionError::OcelNonConforming {
            breed:      input.breed.clone(),
            violations: conformance.violations,
            fitness:    conformance.fitness,
        });
    }

    // -- 5. Temporal conformance -----------------------------------------------
    // check_temporal_conformance verifies that logical_step is strictly
    // monotone across all events in the log.  Out-of-order steps indicate
    // a bug in the breed implementation (e.g., loop that resets the counter).
    check_temporal_conformance(&ocel_log)?;

    // -- 6. Compute BLAKE3 receipt chain --------------------------------------
    // output_hash covers the full BreedOutput, including the ocel_log field.
    // The receipt therefore attests to the process, not just the result.
    // run_id  = blake3(breed || output_hash)
    // replay_pointer = output_hash[..16]
    let output_hash    = blake3_output(&output);
    let run_id         = blake3_run_id(&input.breed, &output_hash);
    let replay_pointer = output_hash[..16].to_string();

    // -- 7. Return ContractResult ----------------------------------------------
    Ok(ContractResult {
        status:          "ok".to_string(),
        breed:           input.breed,
        run_id,
        output_hash,
        replay_pointer,
        options_profile: input.options.map(|o| o.profile).flatten()
                              .unwrap_or_else(|| "default".to_string()),
        output,
    })
}
\end{verbatim}

\noindent
\textbf{Key invariant.}
Because the BLAKE3 \texttt{output\_hash} is computed over a \texttt{BreedOutput}
that contains \texttt{ocel\_log}, the receipt cryptographically binds the
process evidence to the result.  An adversary who wishes to forge a receipt
with a different process trace must break BLAKE3.

\section{Determinism Properties}
\label{app:determinism}

The 13-breed bit-exact double-run harness in
\texttt{tests/breed\_determinism.rs} (16 tests) passes because the following
five design decisions eliminate all sources of non-determinism:

\begin{enumerate}

\item \textbf{No \texttt{HashMap}/\texttt{HashSet} iteration feeding output.}
All internal collections whose iteration order could affect the output are
replaced with \texttt{BTreeMap} (ordered by key) or \texttt{Vec} (sorted
before use).  The Hearsay blackboard, the SOAR preference store, and the
CBR case base all use \texttt{BTreeMap}.  Random-order \texttt{HashMap}
iteration is permitted only for ephemeral scratch structures whose results
are immediately aggregated into a deterministic form (e.g., a sum).

\item \textbf{\texttt{f32::total\_cmp} on clamped \texttt{[0,1]} values.}
All certainty factors, similarity scores, and confidence values are clamped
to \texttt{[0.0,\ 1.0]} before any comparison or sort.  Sorting uses
\texttt{f32::total\_cmp}, which imposes a total order on all finite floats
(including \texttt{-0.0} vs.\ \texttt{+0.0}) and treats \texttt{NaN} as
greater than any finite value.  Because clamping eliminates \texttt{NaN}
and \texttt{+-inf}, the sort is stable and reproducible across runs and
platforms.

\item \textbf{Lexicographic id tiebreaks (smallest key wins on CF tie).}
When two candidates have equal certainty factors after \texttt{total\_cmp},
the tiebreak is lexicographic on the candidate's string identifier
(e.g., rule id, hypothesis label, case id).  The smallest key wins.
This ensures a canonical ordering even when the underlying collection was
populated in insertion order.

\item \textbf{Sorted \texttt{Vec} from \texttt{retrieve\_candidates} (CBR
\texttt{HashSet}$\to$\texttt{Vec} at call site).}
The CBR case base stores retained cases in a \texttt{BTreeMap}, but
\texttt{retrieve\_candidates} previously returned a \texttt{HashSet} of
matching case ids.  The call site now converts this to a \texttt{Vec} and
sorts it lexicographically before scoring.  This was the last remaining
source of non-determinism in the CBR breed; the fix was validated by the
double-run harness.

\item \textbf{\texttt{serde\_json::to\_string} equality for double-run harness.}
The harness serialises the full \texttt{BreedOutput} to JSON using
\texttt{serde\_json::to\_string} (deterministic key order because all maps
are \texttt{BTreeMap} at the serde layer) and asserts byte-exact string
equality between two independent runs with the same input.  This test
surface is stronger than \texttt{PartialEq} on the struct because it also
catches non-determinism introduced by optional fields with default values
or by floating-point formatting differences across serde versions.

\end{enumerate}

\noindent
Together these five properties guarantee that \texttt{wasm4pm-cognition}
satisfies the project's merge-gate determinism requirement:
\emph{same input $\Rightarrow$ bit-exact output}.
